\documentclass[11pt]{article}
\usepackage{preamble}
\renewcommand{\answer}[1]{}

\begin{document}

\ladderheading{AP Statistics \textperiodcentered\ Topic 6.5 (CED 3.6) \textperiodcentered\ B: 2026-12-04 \textperiodcentered\ E: 2027-01-11}{Which Tail, and Probability of What?}

\focusbanner{TODAY'S PROBLEM}

Suppose a city's earlier planning model says $55\%$ of households support expanded evening bus service. A researcher suspects support has \emph{fallen}. From the city's 20{,}000 households, a simple random sample of 200 finds 98 that support expansion, so $\hat p=0.49$. Conditions for a one-sample $z$-test are met.\par\smallskip \textbf{Priya:} ``For $H_0:p=0.55$ versus $H_a:p<0.55$, use the left-tail p-value, about $0.044$. If support is still $55\%$, about $4.4\%$ of random samples of 200 would produce a sample proportion of $0.49$ or lower.''\par \textbf{Noah:} ``Any difference from $0.55$ counts, so double the tail to about $0.088$. That means there is an $8.8\%$ chance that the actual support proportion differs from $0.55$.''\par
\begin{center}
\textbf{Sample counts}\par\smallskip
\begin{tabular}{|l|c|c|c|}
\hline
\textbf{Group} & \textbf{Support} & \textbf{Other} & \textbf{Total} \\ \hline
\textbf{Households} & 98 & 102 & 200 \\ \hline
\end{tabular}
\end{center}
\begin{firsttakebox}{FIRST TAKE --- before the video (5 min)}
Which p-value and interpretation answer whether support has fallen? Choose Priya or Noah and cite one phrase or number.
\workspace{0.9in}
\end{firsttakebox}

\begin{callout}{\IconBook\ A p-VALUE CONDITIONS ON THE NULL (keep on the page)}{calloutgreen}
For a one-sample proportion test,
$z=\dfrac{\hat p-p_0}{\sqrt{p_0(1-p_0)/n}}$.
The p-value is the probability of a sample statistic \textbf{as extreme or more extreme} than the observed one,
in the direction of $H_a$, assuming $H_0$ and the probability model are true. Use the left tail for $H_a:p<p_0$,
the right tail for $H_a:p>p_0$, and both tails for $H_a:p\neq p_0$.
A p-value is a probability about possible sample results under $H_0$; it is not the probability that a hypothesis is true.
\end{callout}

\newpage
\textbf{Finish after the video:}\par\smallskip

\dokbadge{1}~\textbf{(a)} State the null and alternative hypotheses and identify the tail direction required by the research question.\par
\workspace{0.7in}

\dokbadge{2}~\textbf{(b)} Compute the one-sample proportion $z$-statistic and its one-sided p-value. Also compute the two-sided p-value Noah would obtain.\par
\workspace{1.4in}

\dokbadge{3}~\textbf{(c)} Choose the warranted p-value and interpretation. Use $H_a$ and $\hat p=0.49$ to define the tail event and null assumption, then explain what Noah's doubled area answers and what his statement gets wrong.\par
\workspace{2.0in}

\begin{sentenceframebox}\raggedright\textbf{Frame:} Because $H_a:p\blank[0.4in]0.55$, the p-value is the \blank[0.8in] tail beyond $\hat p=\blank[0.6in]$.\end{sentenceframebox}

\begin{sentenceframebox}\raggedright\textbf{Frame:} Assuming \blank[1.8in], a p-value of \blank[0.7in] is the probability that \blank[3.2in], not the probability that \blank[2.0in].\end{sentenceframebox}

\begin{summaryexitbox}{EXIT --- turn this in}
One thing the video changed about my first take: \hrulefill\par
\workspace{0.4in}
\end{summaryexitbox}

\end{document}
